% SPDX-License-Identifier: CC-BY-SA-4.0
% Author: Matthieu Perrin
% Part: <Nom de la partie>
% Section: <Nom de la section>
% Sub-section: <Nom de la sous-section>  % (facultatif, laisser vide si non utilisé)
% Frame: <Titre de la slide>

\begingroup

\begin{frame}{Problème de Correspondance de Post}

  \on[text, top]{
    Soit $\Sigma$ un alphabet contenant au moins deux symboles. 
  }
  
  \on[text, top=4mm, width=1.1\textwidth]{
    \Probleme{pcp}{
      Un ensemble fini de \structure{sortes de dominos} $D \in \mathscr{P}(\Sigma^+ \times \Sigma^+)$.
      \begin{itemize}
      \item Dominos de la forme \VDomino{$\alpha$}{$\beta$}, avec $\alpha$ et $\beta$ des mots sur $\Sigma$
      \end{itemize}
    }{
      Existe-t-il un mot
      $\VDomino{\alpha_1}{\beta_1} \,\cdots\, \VDomino{\alpha_k}{\beta_k} \in D^+$ de dominos\footnote[frame]{On peut utiliser chaque domino autant de fois qu'on veut} tel que :
      \begin{itemize}
      \item $\structure{\alpha_{1} \cdots \alpha_{k}} = \example{\beta_{1} \cdots \beta_{k}}$
      \end{itemize}
    }
  }

  \onExampleBlock[bottom=3mm] {Exemples : les instances suivantes sont-elles positives ?} {
    \begin{itemize}
    \item $\left\{\VDomino{a}{baa}, 
      \VDomino{ab}{aa}, 
      \VDomino{bba}{bb}\right\}
      \uncover<2->{
        \quad\example{\in\textsc{pos}(\textsc{pcp})}
        \quad \text{solution : }
        \VDomino{bba}{bb}\,
        \VDomino{ab}{aa}\,
        \VDomino{bba}{bb}\,
        \VDomino{a}{baa}
      }$
    \item $\left\{\VDomino{a}{ba}, 
      \VDomino{ab}{b}\right\}
      \uncover<2->{\quad\alert{\in\textsc{neg}(\textsc{pcp})}}$
    \end{itemize}
  }
  
\end{frame}

\endgroup
